% Part: proof-theory
% Chapter: natural-deduction
% Section: grafting

\documentclass[../../../include/open-logic-section]{subfiles}

\begin{document}

\olfileid[id]{pt}{nat}{gra}

\olsection{Pencangkokan \usetoken{P}{proof}}

!!^{proof} dalam \Log{N1c} dan~\Log{N1i} merupakan !!{proof} atas
!!{formula}~$!A$ dari asumsi-asumsi~$!B$. Andaikan $!B$ sendiri memiliki
!!a{proof}. Bagaimana kita menggabungkan !!{proof} atas~$!B$ dengan
!!{proof} atas $!A$ dari~$!B$ untuk memperoleh !!a{proof} atas~$!A$ yang
tidak menggunakan asumsi~$!B$?

Salah satu pilihan ialah menghapus asumsi~$!B$ dari !!{proof} atas~$!A$
dengan menerapkan~\Intro{\lif}. Kemudian kita dapat menggabungkan
!!{proof} yang dihasilkan atas~$!B \lif !A$ dengan !!{proof} atas~$!B$
untuk memperoleh
\[
\AxiomC{$\Discharge{!B}{x}$}
\DeduceC{$!A$}
\DischargeRule{\Intro{\lif}}{x}
\UnaryInfC{$!B \lif !A$}
\AxiomC{}
\DeduceC{$!B$}
\RightLabel{\Elim{\lif}}
\BinaryInfC{$!A$}
\DisplayProof
\]
Cara ini langsung, tetapi agak janggal: mengapa mengintroduksi $\lif$ hanya
untuk segera mengeliminasinya? Semestinya kita dapat menghindari
``jalan memutar'' semacam ini melalui $!B \lif !A$. (Fakta bahwa jalan
memutar itu selalu dapat dihindari sebenarnya merupakan isi teorema
normalisasi.)

Dalam \Log{N1c} dan~\Log{N1i}, kita dapat menghindarinya dengan cukup mudah
dalam kasus ini: kita cukup mengganti asumsi-asumsi~$!B^x$ dengan seluruh
!!{proof} atas~$!B$. ``Pencangkokan'' !!{proof} pada asumsi lain ini
cukup berguna, sehingga kita mencatatnya sebagai definisi.

\begin{defn}
Jika $\delta_1$ adalah \Log{N1c}-!!{proof} (atau
\Log{N1i}-!!{proof}) atas~$!A$ dari asumsi terbuka~$!B^x$, dan
$\delta_2$ adalah \Log{N1c}-!!{proof} (\Log{N1i}-!!{proof})
atas~$!B$, maka $\Subst{\delta_1}{\delta_2}{!B^x}$ adalah hasil
penggantian setiap kemunculan $!B^x$ yang belum dilepaskan
dalam~$\delta_1$ dengan~$\delta_2$.
\end{defn}

Asalkan $\delta_1$ dan $\delta_2$ tidak memiliki !!{formula} asumsi yang berbeda
tetapi diberi label yang sama, dan tidak ada asumsi terbuka
dalam~$\delta_2$ yang memuat variabel eigen dari~$\delta_1$, hasil
$\Subst{\delta_1}{\delta_2}{!B^x}$ merupakan !!{proof} yang benar, dan
asumsi-asumsi terbukanya ialah asumsi-asumsi terbuka $\delta_1$ bersama
dengan asumsi-asumsi terbuka~$\delta_2$. Ketika mencangkokkan !!{proof},
syarat-syarat ini biasanya terpenuhi. Syarat pertama terpenuhi setiap kali
$\delta_1$ dan $\delta_2$ merupakan subbukti dari !!{proof} yang lebih
besar. Jika syarat kedua tidak terpenuhi, kita dapat mengganti nama
variabel-variabel eigen dalam~$\delta_1$ agar syarat itu terpenuhi.

Definisi yang bersesuaian berlaku bagi \Log{N2c} dan~\Log{N2i}.

\begin{defn}
Jika $\delta_1$ adalah \Log{N2c}-!!{proof} (atau
\Log{N2i}-!!{proof}) atas~$x:!B, \Gamma_1 \Sequent !A$, dan
$\delta_2$ adalah \Log{N2c}-!!{proof} (\Log{N2i}-!!{proof})
atas~$\Gamma_2 \Sequent !B$, maka
$\Subst{\delta_1}{\delta_2}{x:!B}$ adalah hasil penggantian setiap
aksioma $x:!B \Sequent !B$ dalam~$\delta_1$ dengan~$\delta_2$, serta
penambahan $\Gamma_2$ pada konteks setiap sekuen di bawah aksioma-aksioma
tersebut.
\end{defn}

\begin{prop}
Jika $\delta_1 \Proves x:!B, \Gamma_1 \Sequent
!A$ dan $\delta_2 \Proves \Gamma_2 \Sequent !B$, maka
$\Subst{\delta_1}{\delta_2}{x:!B} \Proves \Gamma_1, \Gamma_2 \Sequent
!A$, asalkan tidak ada variabel eigen dari~$\delta_1$ yang muncul
dalam~$\Gamma_2$.
\end{prop}

\begin{proof}
  Melalui induksi pada $\pheight{\delta_1}$.
\end{proof}

\end{document}
